\rtmtight
\begin{tblr}{width=\textwidth,colspec={|Q[c,m,2.1cm]|X[l,m]|},hlines,vlines,hspan=minimal,rowsep=1.5pt,colsep=3pt,row{1}={bg=headgray,font=\bfseries}}
\SetCell{c,m}\hfil\logo\hfil & \textbf{POLITEKNIK NEGERI MALANG}\par\textbf{JURUSAN TEKNOLOGI INFORMASI}\par\textbf{PROGRAM STUDI: D4 TEKNIK INFORMATIKA} \\
\SetCell[c=2]{l} \textbf{RENCANA TUGAS MAHASISWA (RTM) DAN RUBRIK PENILAIAN} & \\
Nama Mata Kuliah & Analisis dan Desain Berorientasi Objek \\
Kode Mata Kuliah & RTI254002 \quad \textbf{sks / Jam perminggu:} 2 SKS/4 jam \quad \textbf{Semester:} 4 \\
Dosen Pengampu & Tim Pengajar Program Studi D4 TI \\
\end{tblr}
\assessmentblock{Tugas Artefak Analisis dan Desain ADBO}{Case Method dan analisis kasus}{SCPMK0211-02901, SCPMK0602-02902, SCPMK1008-02903}{Mahasiswa menghasilkan artefak ADBO secara bertahap setiap minggu, mencakup diagram UML, spesifikasi use case, dan spesifikasi arsitektur teknis.}{Menganalisis kebutuhan kasus yang diberikan, membuat diagram UML yang diminta, mendokumentasikan keputusan, dan mengumpulkan evidence.}{Diagram UML, spesifikasi use case, dokumen SKPL, laporan, dan/atau bahan presentasi.}{Kualitas artefak analisis UML & 12 & Use case, class diagram, sequence, activity, dan state diagram disusun dengan benar, notasi tepat, dan konsisten dengan kebutuhan sistem. & Artefak UML sebagian besar benar tetapi beberapa diagram belum konsisten atau notasi masih ada kesalahan minor. & Artefak UML tidak mencerminkan analisis kebutuhan yang benar atau banyak kesalahan notasi. \\ \\
Ketepatan pemodelan struktural UML lanjutan & 10 & Advanced class diagram, state machine, dan deployment diagram dipilih dengan tepat sesuai konteks, diimplementasikan dengan benar, dan justifikasi pemodelan terdokumentasi jelas. & Pemodelan struktural cukup tepat tetapi implementasi atau justifikasinya belum lengkap. & Pemodelan struktural tidak tepat, tidak diterapkan sesuai konteks, atau tidak ada justifikasi. \\ \\
Kelengkapan spesifikasi teknis ADBO & 8 & SRS, diagram arsitektur, component dan deployment diagram terintegrasi menjadi spesifikasi yang siap diimplementasikan. & Spesifikasi sebagian besar lengkap tetapi beberapa artefak belum terintegrasi atau konsisten. & Spesifikasi tidak lengkap atau artefak tidak saling terhubung. \\}

\assessmentblock{Kuis Konsep UML}{Kuis}{SCPMK0211-02901, SCPMK0602-02902}{Kuis tertulis untuk mengukur pemahaman domain model dan use case modeling (Kuis 1) serta sequence diagram dan class diagram (Kuis 2).}{Menjawab soal pilihan ganda dan/atau uraian terkait konsep UML, notasi diagram, use case specification, sequence diagram, dan class diagram.}{Jawaban tertulis kuis.}{Penguasaan konsep UML dan analisis OO & 3 & Use case, class diagram, dan diagram analisis dijelaskan dan diterapkan dengan benar pada kasus yang diberikan. & Sebagian besar konsep UML dipahami tetapi penerapannya pada kasus belum sepenuhnya tepat. & Konsep UML tidak dipahami atau banyak kesalahan dalam penerapan notasi. \\ \\
Kemampuan membaca dan menginterpretasi diagram & 4 & Diagram UML dibaca dan diinterpretasikan dengan benar, pertanyaan tentang diagram dijawab secara akurat. & Sebagian besar diagram dapat dibaca dengan benar tetapi beberapa interpretasi masih keliru. & Tidak dapat menginterpretasi diagram UML dengan benar. \\ \\
Penguasaan pemodelan struktural UML & 3 & Sequence diagram dan class diagram (relationships, visibility, attribute, operation) dijelaskan dengan benar dan contoh penerapannya pada kasus tepat. & Sebagian besar konsep dipahami tetapi contoh atau penerapannya belum lengkap. & Konsep tidak dipahami atau contoh penerapannya keliru. \\}

\assessmentblock{UTS Artefak Analisis ADBO}{UTS berbasis artefak}{SCPMK0211-02901, SCPMK0602-02902}{Evaluasi tengah semester berupa presentasi progress SKPL untuk mengukur kemampuan menghasilkan artefak analisis ADBO (domain model, use case, robustness analysis, dan activity diagram).}{Menganalisis kasus baru yang diberikan, menghasilkan domain model, use case specification, robustness analysis, dan activity diagram.}{Dokumen artefak ADBO: domain model, use case specification, robustness analysis, dan activity diagram.}{Kelengkapan dan ketepatan diagram analisis & 5 & Domain model, use case, robustness analysis, dan activity diagram disusun lengkap, notasi benar, dan konsisten antara satu diagram dengan lainnya. & Diagram utama ada tetapi beberapa notasi atau konsistensi antara diagram belum tepat. & Diagram tidak lengkap atau notasi banyak kesalahan. \\ \\
Kualitas use case specification dan skenario & 4 & Use case specification ditulis lengkap dengan skenario utama, alternatif, dan kondisi pengecualian yang detail dan konsisten. & Spesifikasi ada tetapi skenario alternatif atau detail kondisi pengecualian belum lengkap. & Tidak ada spesifikasi use case yang memadai atau skenario utama tidak jelas. \\ \\
Kualitas robustness analysis & 4 & Robustness analysis (boundary, control, entity object) diidentifikasi dengan tepat dan konsisten dengan use case yang dianalisis. & Robustness analysis ada tetapi identifikasi elemen belum lengkap atau kurang konsisten. & Robustness analysis tidak dilakukan atau tidak sesuai dengan use case yang dianalisis. \\}

\assessmentblock{PBL Desain Sistem ADBO Terpadu}{Project Based Learning}{SCPMK0602-02902, SCPMK1008-02903}{Mahasiswa mengintegrasikan seluruh artefak ADBO yang telah dikerjakan menjadi dokumen SKPL/desain sistem lengkap dan mempresentasikan hasilnya kepada dosen dan rekan mahasiswa.}{Mengintegrasikan SRS, model domain, diagram UML (class diagram lanjutan, state machine, component, dan deployment diagram) menjadi satu dokumen SKPL yang kohesif, lalu mempresentasikannya.}{Dokumen SKPL/desain sistem ADBO terpadu, diagram-diagram, dan bahan presentasi.}{Kualitas pemodelan struktural dan arsitektur UML & 9 & Class diagram lanjutan, state machine, component, dan deployment diagram diterapkan konsisten, saling terhubung, dan dibenarkan dengan argumen teknis yang kuat. & Pemodelan struktural cukup baik tetapi beberapa justifikasi atau konsistensi antarkomponen belum kuat. & Pemodelan struktural tidak konsisten dengan domain model atau tidak diterapkan dengan tepat. \\ \\
Kelengkapan artefak desain terpadu & 8 & SRS, model domain, diagram arsitektur, component, deployment, dan spesifikasi teknis hadir, terintegrasi, dan saling konsisten. & Sebagian besar artefak ada tetapi integrasi atau konsistensi antartefak masih kurang. & Artefak tidak terintegrasi atau spesifikasi tidak siap diimplementasikan. \\ \\
Presentasi dan argumentasi keputusan desain & 5 & Keputusan desain disampaikan dengan jelas, didukung argumen teknis, dan dapat dipertahankan dalam sesi tanya jawab. & Presentasi cukup jelas tetapi beberapa keputusan tidak dapat dipertahankan dengan baik. & Presentasi tidak menjelaskan keputusan desain atau argumen tidak relevan. \\}

\assessmentblock{UAS Review Dokumen Desain Sistem}{UAS berbasis presentasi dan defense}{SCPMK1008-02903}{Mahasiswa mempresentasikan dokumen desain sistem final secara menyeluruh dan mempertahankan setiap keputusan desain kepada dosen penguji.}{Menyiapkan presentasi komprehensif dokumen desain sistem, mendemonstrasikan konsistensi antarartefak, dan menjawab pertanyaan penguji terkait keputusan desain.}{Dokumen desain sistem final, slide presentasi, dan/atau demonstrasi diagram.}{Kelengkapan dan konsistensi spesifikasi teknis & 10 & Semua artefak ADBO (SRS, model domain, diagram arsitektur, component, deployment) hadir, konsisten, dan implementable. & Sebagian besar spesifikasi lengkap tetapi beberapa artefak atau konsistensi antarartefak belum terpenuhi. & Spesifikasi tidak lengkap atau banyak inkonsistensi yang menghambat implementasi. \\ \\
Kualitas konsistensi pemodelan struktural UML & 9 & Class diagram lanjutan, state machine, dan deployment diagram diterapkan konsisten, justifikasi pemilihan kuat, dan dapat ditelusuri dalam seluruh dokumen SKPL. & Pemodelan struktural diterapkan tetapi justifikasi atau konsistensinya belum menyeluruh di semua bagian dokumen. & Pemodelan struktural tidak diterapkan dengan benar atau justifikasi tidak ada. \\ \\
Defense dan review dokumen desain & 6 & Mahasiswa mampu menjelaskan dan mempertahankan setiap keputusan desain dengan evidence dari dokumen dan argumen teknis. & Defense cukup baik tetapi beberapa keputusan tidak dapat dipertahankan atau dokumen memiliki inkonsistensi. & Tidak dapat menjelaskan keputusan desain atau dokumen banyak kesalahan yang tidak dapat dipertahankan. \\}

\begin{tblr}{width=\textwidth,colspec={|Q[l,m,3.2cm]|X[l,m]|},hlines,vlines,hspan=minimal,row{1-3}={bg=headgray},rowsep=1.5pt,colsep=3pt}
Jadwal Pelaksanaan: & Minggu 1--17 sesuai RPS. Setiap evaluasi memiliki RTM dan rubrik multi-kriteria. \\
Lain-Lain Yang Diperlukan: & LMS, repositori artefak desain, perangkat lunak diagram UML, dan perangkat presentasi. \\
Pustaka: & Mengacu pustaka utama dan pendukung pada bagian RPS. \\
\end{tblr}
